% GENERATED by scripts/make_waste_completeness_table.py - do not
% hand-edit; a test asserts this file matches the engine's completeness table.
% ===========================================================================
\section{Supplementary Table: the IPCC sector-4 completeness matrix}
\label{app:completeness}
% ===========================================================================

Every IPCC waste category, with its notation key (E estimated, NE not
estimated, IE included elsewhere) and the justification the engine carries
for every non-E entry. Generated from the same completeness table the
national inventory publishes, which is itself synced bidirectionally to the
gap register by test -- an NE row without a register entry, or a register
entry without an NE row, fails the build.

\begin{table*}[htbp]
\centering
\scriptsize
\begin{tabular}{@{}llclp{0.43\linewidth}@{}}
\toprule
\textbf{Code} & \textbf{Category} & \textbf{Key} & \textbf{Gases} & \textbf{Justification (abridged)} \\
\midrule
4.A & Solid waste disposal - municipal & E & CH4 & IPCC Vol 5 Ch 3 First Order Decay over a 1960-2024 deposit series, driven by the COLLECTED complement of the open-burning stream, so the two categories partition one generation figure and cannot double-count a tonne. Degradable organic carbon is the biogenic\ldots \\[2pt]
4.A & Solid waste disposal - industrial waste & NE & CH4 & The estimated category is municipal solid waste only. Industrial solid waste reaching disposal sites has no national tonnage here; the provincial EPAs' industrial waste returns are the closing source. See waste-industrial. \\[2pt]
4.A & Solid waste disposal - construction waste & NE & CH4 & Construction and demolition waste is largely inert (near-zero degradable carbon), so its CH4 potential is small by composition - but its tonnage is unquantified and it is a distinct absence from industrial waste, with a different closing source: the Urban Uni\ldots \\[2pt]
4.A & Solid waste disposal - CO2 from landfill gas & IE & CO2 & Biogenic, and unlike open burning IPCC does not request it as a memo item: that carbon is tracked as a stock change in AFOLU. \\[2pt]
4.B & Biological treatment of solid waste (composting, anaerobic\ldots & NE & CH4/N2O & Municipal composting operates at city scale - Lahore's composting facility is the largest, with smaller plants elsewhere - but no national throughput series has been compiled here. Believed small against 4.A, and the omission biases 4.A UP rather than the tot\ldots \\[2pt]
4.C.1 & Waste incineration & NE & CO2/CH4/N2O & No municipal-scale incineration operates in Pakistan, but hospital and clinical-waste incinerators do; their throughput has not been obtained here. Bounded far below 4.C.2, whose 3.83 Mt of open burning dwarfs any clinical stream. See waste-incineration. \\[2pt]
4.C.2 & Open burning of waste & E & CO2/CH4/N2O & IPCC Eq 5.7 over the UNCOLLECTED stream x a 0.60 burned fraction. Only fossil CO2 counts; biogenic CO2 is a memo item. \\[2pt]
4.D.1 & Domestic wastewater & E & CH4/N2O & IPCC Vol 5 Ch 6 on the district population panel, stratified by the WHO/UNICEF JMP sanitation split - sewer 62.6\% urban against 8.8\% rural, which is why the stratification matters. THE JUDGEMENT THAT DECIDES THIS CATEGORY is the sewer MCF: IPCC gives stagna\ldots \\[2pt]
4.D & Industrial wastewater & NE & CH4 & Needs industrial output by sector and wastewater organic loading, neither of which is available; blocked with 1.A.2 on the same HDIP/industry data. \\[2pt]
\bottomrule
\end{tabular}
\caption{The waste-sector completeness matrix, generated from the engine
(\texttt{outputs/completeness.csv}, sector-4 rows). Full justifications are
in the generated national table; each NE row names its gap-register entry.}
\label{tab:completeness}
\end{table*}
